\section{Computational study}
\label{sec:experiments}

The experiments pursue three questions. (Q1) Is the enumeration algorithm of
Theorem~\ref{thm:xp} practical in its predicted regime, and where does that
regime end. (Q2) How does it compare with the standard mixed-integer approach.
(Q3) Does the covering bound $k^\ast$ track instance difficulty beyond the
algorithm that motivated it, and through which mechanism. The remainder of the section is organised so that
the reader can locate each answer directly. Section~\ref{sec:exp-setup}
describes the codes, the instance families and the machines.
Section~\ref{sec:exp-q1} answers Q1 on the family whose covering bound is
controlled by construction, Section~\ref{sec:exp-q2} addresses Q2 against
the integer programming baseline, and Section~\ref{sec:exp-q3} addresses
Q3 by
confronting $k^\ast$ with the effort of a solver that makes no use of it.

\subsection{Setup}
\label{sec:exp-setup}

The study compares three exact codes. The first (\textsc{Xp}) is a C++20
implementation of Algorithm~\ref{alg:xpbb}, with the covering pruning, the
admissible cost bound and the dominance rule of Section~\ref{sec:algorithms},
and a deadline-aware enumerator. The baseline (\textsc{Mip}) is the compact
model \eqref{eq:obj}--\eqref{eq:dom} solved by Gurobi. A third code
(\textsc{Bd}) solves the same model by a branch-and-Benders-cut in which the
per-product transportation subproblems are separated by the disaggregated cuts
of Proposition~\ref{prop:duals}. All three read the same instance files.
The \textsc{Xp} implementation was cross-validated against an independent
brute-force enumeration of all designs with an exact per-product min-cost-flow
oracle, matching it on $180$ small seeded instances and on $1{,}381$ exact
comparisons of optimal value, feasibility and covering bound across a full sweep
of the cardinality, with no disagreement. The verification battery in the
accompanying code performs on the order of nine hundred thousand
individual checks, spanning design enumerations, closed-form comparisons
and LP cross-checks, with no failure. More than $560{,}000$ of the
checks are exhaustive design enumerations validating the numeric-layer
constructions alone, roughly ten thousand are equivalence
verifications, two hundred are independent LP cross-checks, $1{,}381$
are cross-solver comparisons, and the remainder are per-construction
parameter sweeps.

The study uses two instance families, the first with a main grid and a
boundary extension, the second the benchmark. The first is a family with controlled
covering bound, built by a deterministic generator seeded by the parameter
tuple, in which, on each side, one facility fewer than the per-side target
receives capacities that jointly fall one unit short of the binding
product's demand, a choice that pins the per-side covering count to the
target. The main
grid has target per-side count in $\{2,3,4,5,6\}$, so $k^\ast\in\{4,6,8,10,12\}$,
factory and depot counts in $\{20,40\}$, customers in $\{50,100\}$, products in
$\{5,10\}$ and five seeds, for $400$ instances, and a boundary extension on
smaller graphs, with per-side target in $\{2,4,6,8\}$ so
$k^\ast\in\{4,8,12,16\}$ and $n=|I|+|J|$ over $\{20,24,28,30,40\}$, probes where
the regime ends, adding $60$ instances for $460$ runs in all. The
generator records the recomputed bound alongside each instance, and on
all $460$ instances the recomputed $k^\ast$ equals the target exactly.
The second family is the classical
benchmark of \citet{MauriEtAl2021}, its $100$ instances split into four size
groups $|I|$-$|J|$-$|K|$ of $50$-$100$-$200$ and $100$-$200$-$400$ with
$|L|\in\{5,10\}$, for which no optimal value was proven in \citet{MauriEtAl2021} nor, to
our knowledge, since. On these instances the customer-aggregation rule of
Proposition~\ref{prop:aggregation} never applies, since in every one of
the $100$ instances no two customers share their full cost column and
hence $\tau=|K|$. Capacity capping leaves every benchmark capacity
unchanged, since no capacity exceeds its product's total demand,
customer aggregation never applies as just noted, and the covering
bound is thus the only signal the paper's toolkit produces on these
instances. One by-product of
parsing them is used repeatedly
below and is recorded here. Every
datum is integral, so the optimal value is integral by
Proposition~\ref{prop:oracle}, and a solver that reports optimality with
an absolute gap below one therefore determines the exact optimum,
subject to the floating-point tolerances under which the solver
certifies its bound, a caveat Section~\ref{sec:exp-q2} makes explicit
for the three closed runs. The
terminal gap of a run is the relative optimality gap the solver reports
when it stops, the difference between incumbent and bound divided by the
incumbent, quoted in percent throughout.

The enumeration runs were executed under a fixed $60$~s limit. The
\textsc{Xp} binary is single-threaded C++20, compiled with g++~13.3.0 at
optimisation level O2, and ran on a Linux virtual machine with an Intel
Xeon processor at $2.1$~GHz.
The integer programming runs use Gurobi $13.0.2$ on an Intel Core i9-14900HX
with $24$ cores under Windows, with the controlled family limited to $600$~s
and the benchmark to $3{,}600$~s. Gurobi ran with default tolerances,
automatic thread count and a fixed seed recorded with every run, and on
each run the harness recomputed the incumbent's cost independently from
its flows, with agreement on all $900$ runs. The enumeration and the
integer programming codes ran on different machines, so times are
compared only within a code and never across codes, and the
solved-versus-censored contrasts that the section draws are insensitive
to the hardware difference. The asymmetry of the limits favours the
baseline by
design. The enumeration limit maps where that algorithm stops being fast, the
larger Gurobi budget gives the baseline its best chance of producing
reference optima, and a censored run enters a speed comparison only
through the lower bound its limit certifies. The time of a censored run
is recorded at its limit. Every completed enumeration run is
deterministic, and each instance was run once. The software archive with
the codes, the instance generator and manifests, the raw result files and
one verification script per hardness construction accompanies the
submission as supplementary material and is released at
\url{https://github.com/pehgonzalez/mp-tscflp-parameterized}.

\subsection{The regime of the enumeration algorithm (Q1)}
\label{sec:exp-q1}

The controlled family was built to probe the boundary that
Theorem~\ref{thm:xp} predicts, and the outcome is unambiguous. The main grid, whose
graphs all have $n=|I|+|J|\ge 40$, lies entirely outside the $60$~s regime, so
all $400$ instances are censored. This is the result rather than a failure of
protocol. At $n\ge 40$ the search processes a median of roughly
$2{,}700$ nodes before the limit, while a completed search at $n=24$
already needs a median of roughly $50{,}000$, and the cost of each node
grows with $|I|$, $|J|$, $|K|$ and $|L|$ through the oracle. The raw
per-run records in the archive carry node, pruning and oracle
counters. The boundary extension locates the frontier, and
Table~\ref{tab:q1boundary} reports it in full. Each boundary cell holds
three seeded runs and the medians are taken over solved runs only, so the
times are indicative and the visible values at $n=24$ are biased downward
by that conditioning. At $n=20$ every instance
solves, at $n=24$ seven of twelve solve, and from $n=28$ upward none solves
within the limit, so the boundary observed for this implementation
under this protocol, in plain optimisation mode,
sits at $n\approx 24$. The same table shows the covering-bound effect at
fixed size, with the median time at $n=20$ falling from $13.7$~s at
$k^\ast=4$ to $0.6$~s at $k^\ast=16$. Figure~\ref{fig:q1} shows both
readings graphically, the boundary in panel (a) and the covering-bound
effect in panel (b).

\begin{table}[H]
\centering
\caption{The boundary extension of Q1 under the $60$~s limit. Each cell
reports solved over attempted instances at the given $n=|I|+|J|$ and
covering-bound target $k^\ast$, with the median time in seconds of the
solved runs in parentheses. Solved runs vanish beyond $n=24$, and in the
solvable rows a larger bound solves faster. Starred cells have fewer
than three solved runs and their medians are individual observations.}
\label{tab:q1boundary}
\begin{tabular}{lcccc}
\hline
$n$ & $k^\ast=4$ & $k^\ast=8$ & $k^\ast=12$ & $k^\ast=16$ \\
\hline
$20$ & 3/3 (13.7) & 3/3 (4.0) & 3/3 (2.3) & 3/3 (0.6) \\
$24$ & 1/3 (58.7)$^{\star}$ & 1/3 (57.2)$^{\star}$ & 2/3 (36.2)$^{\star}$ & 3/3 (15.0) \\
$28$ & 0/3 & 0/3 & 0/3 & 0/3 \\
$30$ & 0/3 & 0/3 & 0/3 & 0/3 \\
$40$ & 0/3 & 0/3 & 0/3 & 0/3 \\
\hline
\end{tabular}
\end{table}


\begin{figure}[H]
\centering
\includegraphics[width=\textwidth]{fig_q1_regime.pdf}
\caption{The enumeration algorithm on the controlled family, all $460$ runs
under the $60$~s limit. Panel (a) plots time against $n=|I|+|J|$. Solved runs
are filled and coloured by the covering bound $k^\ast$, and they occur only in
the shaded band up to $n\approx 24$. Every run beyond it is censored, shown as
grey open circles on the limit line, including the whole $400$-instance main
grid, which lies at $n\ge 40$. Panel (b) gives the median solve time against $k^\ast$ at
$n=20$ and $n=24$, where instances with a larger bound solve faster,
with medians taken over solved runs only. Open markers in panel (b) are
medians over fewer than three solved runs. Lower is better in both
panels, and the censored band beyond $n\approx 24$ is the observed
border of the regime that Theorem~\ref{thm:xp} delimits.}
\label{fig:q1}
\end{figure}

Within the solvable range the covering bound acts in a direction that
requires explanation. Time and node counts decrease as $k^\ast$ grows,
and the solved count at $n=24$ in Table~\ref{tab:q1boundary} rises from
one in three at $k^\ast=4$ to three in three at $k^\ast=16$. The
covering pruning of Lemma~\ref{lem:cover} is the mechanism compatible
with this pattern. In plain optimisation the
cardinality never binds, so $k^\ast$ can act on the search only through
that pruning, and tight capacities let it discard subtrees
that loose capacities leave standing. The generator, however, varies
capacities to move the bound and therefore moves other instance
features with it, so an ablation isolating the pruning is left for
future work. The mirror pattern expected for the
cardinality sweep, in which the exponent $n^k$ grows with the target count, is
left for a dedicated run.

\subsection{Comparison with the mixed-integer approach (Q2)}
\label{sec:exp-q2}

Q2 places the enumeration and the integer programming codes side by
side, each under its own protocol, and the observed outcome is an
empirical result consistent with the worst-case limitation of support
enumeration rather than a consequence of it. Outside the
small-$n$ regime the weight-independent enumeration bound gives \textsc{Xp} no
way to prune below $n^{k}$, whereas the compact model carries the linear
structure a modern solver exploits, so \textsc{Mip} is expected to close quickly the very controlled-family
instances on which \textsc{Xp} exhausts its $60$~s budget. The campaign
realises that expectation, and Table~\ref{tab:q2codes} summarises it.
\textsc{Mip} closes all $400$ controlled instances with a median solving
time of $0.55$~s, among them every instance the enumeration had left
censored at $60$~s.

\begin{table}[H]
\centering
\caption{The three exact codes on the $400$-instance main grid of the
controlled family, whose graphs all have $n\ge 40$. Solved means proven
optimal within the per-code limit. The median time is over solved runs
only, in seconds. \textsc{Mip} closes every instance the enumeration
leaves censored, the empirical face of the regime border of Q1. The
enumeration ran on a different machine from the two integer programming
codes, which shared one machine, so times compare only within a code
and between \textsc{Mip} and \textsc{Bd}, and the cross-code reading
is the solved-versus-censored contrast.}
\label{tab:q2codes}
\begin{tabular}{lccc}
\hline
code & limit (s) & solved & median time (s) \\
\hline
\textsc{Xp} & $60$ & $0/400$ & -- \\
\textsc{Mip} & $600$ & $400/400$ & 0.55 \\
\textsc{Bd} & $600$ & $8/400$ & 395.57 \\
\hline
\end{tabular}
\end{table}
 The value of
Theorem~\ref{thm:xp} is thus classificatory rather than competitive, and
the observed dominance is no defeat of the algorithm but an
empirical face consistent with the same hardness the theory maps. What the exact codes
earn is not speed but proof, since by the integrality remark of
Section~\ref{sec:exp-setup} a solver that closes an instance settles a
value the benchmark of \citet{MauriEtAl2021} had left to heuristics. On
the benchmark itself the $3{,}600$~s budget proves optimality on three of
the $100$ instances, all in the smaller size group with five products,
values the original study left unproven. The three instances are
PSC1-C1-50-5, PSC4-C1-50-5 and PSC5-C1-50-5, closed at optimal values
$1{,}923{,}801$, $1{,}909{,}659$ and $1{,}895{,}151$ in $2{,}001$, $642$
and $883$ seconds. Each value lies strictly below the best solution
reported for the same instance by \citet{MauriEtAl2021}, respectively
$1{,}923{,}813$, $1{,}910{,}011$ and $1{,}896{,}699$, so the campaign
closes the three instances and improves their best-known solutions at the
same time. On these three runs the terminal bound equalled the incumbent
at the recorded precision, and we report the optima subject to the
solver's default floating-point certificates, with the incumbent costs
confirmed by the independent recomputation noted above.
The branch-and-Benders-cut \textsc{Bd} enters here not as a faster
solver, for it closes $8$ of the $400$ controlled instances within the
same budget, but as a device for cross-validation and for the structural
reading of Proposition~\ref{prop:duals}, and the codes are compared only
on optimal values they jointly certify.

\subsection{The covering bound and the root relaxation (Q3)}
\label{sec:exp-q3}

The last question asks whether $k^\ast$ tracks difficulty beyond the
algorithm that motivated it, and through which mechanism, so we
confront the covering bound with the integer
solver on the benchmark that the original study left without a single proven
optimum. The bound is available
without solving anything, computed by one sort and one prefix scan,
and it varies widely across the benchmark, as Figure~\ref{fig:q3} shows. Table~\ref{tab:q3groups} reports the campaign by
size group, the range of the bound, the median terminal gap, the proven
optima and the within-group correlations with their Holm-adjusted
$p$-values. Within every size group the
bound spreads over a factor of four or more, and this spread at fixed
size removes the most evident confounder, instance size, while other
instance features still vary within a group, so the bound may in part
proxy properties the design does not separate.

\begin{table}[H]
\centering
\caption{The benchmark campaign by size group, $3{,}600$~s per instance.
Columns give the group $|I|$-$|L|$, the instance count, the minimum,
median and maximum of the covering bound $k^\ast$, the median terminal
gap in percent, the number of proven optima, the within-group Spearman
correlation between $k^\ast$ and the terminal gap, and its
Holm-adjusted two-sided permutation $p$-value. Correlations carry
percentile bootstrap $95\%$ confidence intervals in brackets. Each
interval describes one group alone, and the within-group column is the
sharper instrument, the pooled value mixing the group effect with size.
The pooled correlation over all $100$ instances is $-0.28$,
Holm-adjusted $p=0.013$.}
\label{tab:q3groups}
\small
\setlength{\tabcolsep}{4pt}
\begin{tabular}{lcccccc}
\hline
group & $\#$ & $k^\ast$ min/med/max & gap (\%) & opt.\ &
$\rho$ [$95\%$ CI] & $p$ \\
\hline
$50$-$5$ & $25$ & $8$/$17$/$36$ & $0.4$ & $3$ & $-0.59$ [$-0.81$, $-0.22$] & $0.006$ \\
$50$-$10$ & $25$ & $8$/$17$/$35$ & $0.4$ & $0$ & $-0.39$ [$-0.69$, $+0.09$] & $0.102$ \\
$100$-$5$ & $25$ & $14$/$32$/$67$ & $1.1$ & $0$ & $-0.70$ [$-0.86$, $-0.45$] & $0.001$ \\
$100$-$10$ & $25$ & $14$/$33$/$69$ & $4.2$ & $0$ & $-0.40$ [$-0.75$, $+0.12$] & $0.102$ \\
\hline
\end{tabular}
\end{table}


\begin{figure}[H]
\centering
\includegraphics[width=0.8\textwidth]{fig_q3_kstar.pdf}
\caption{The covering bound $k^\ast$ across the four benchmark groups. Each
point is one instance and the bar marks the group median. The bound spreads by
a factor of more than four inside every group, varying across instances of
equal size, so its variation is not a proxy for instance size.}
\label{fig:q3}
\end{figure}

The completed campaign lets the correlation between $k^\ast$ and the
effort \textsc{Mip} spends on the $100$ benchmark instances be read
directly, with a within-group control that separates the effect of the
bound from the effect of size. Pooled over the benchmark, the covering
bound and the terminal gap are rank-correlated with Spearman
$\rho=-0.28$. Ninety-seven of the hundred runs end at the limit, so
solving time carries no usable signal on this benchmark and the analysis
reads the terminal gap alone. The terminal gap blends primal and dual
progress and depends on the objective scale and on the solver
tolerances, so it is an imperfect proxy of computational effort.
Within the four size groups the association strengthens, the four
coefficients being negative throughout and ranging from
$-0.39$ to $-0.70$ with descriptive mean $-0.52$, as the $\rho$ column
of Table~\ref{tab:q3groups}
shows. Significance was assessed by a two-sided permutation test on the
Spearman statistic. The test is Monte Carlo with $20{,}000$
permutations under a fixed seed, the $p$-value uses the add-one
convention so it cannot vanish, ties including zero gaps are handled by
mid-ranks, and every reported $p$-value is Holm-adjusted over the five
gap correlations. The pooled correlation and the two five-product
groups survive the correction, with adjusted $p$-values of $0.013$,
$0.006$ and $0.001$, while the two ten-product groups do not, and the
$p$ column of Table~\ref{tab:q3groups} reports the adjusted value of
every group. The direction is negative in all four groups, established
where the test passes and inconclusive where it does not. The bootstrap
intervals of Table~\ref{tab:q3groups} tell the same story, excluding
zero exactly in the groups that survive the correction.
Figure~\ref{fig:q3scatter} plots the hundred instances behind these
coefficients, and the negative drift is visible inside each group.
The sign is the substantive finding. A naive reading
of the bound would predict the opposite, a small $k^\ast$ meaning few
facilities to open and hence an easy instance, but on this benchmark
larger bounds go with smaller terminal gaps, the same inverse direction
the enumeration exhibits in Q1, where tight capacities give the covering
pruning its greatest effect.

The stored solver logs admit a mediation analysis that locates the
mechanism, extracted by a pinned script with no further runs. The
terminal gap on this benchmark is essentially the root relaxation gap,
the two rank-correlating between $+0.89$ and $+1.00$ within groups and
at $+0.97$ pooled. The covering bound correlates with the root gap
between $-0.42$ and $-0.73$ in every group ($-0.36$ pooled), and once
the root gap is controlled the partial rank correlation of the bound
with the terminal gap collapses, ranging from $-0.04$ to $+0.52$ across
groups with no stable sign ($+0.29$ pooled). The bound moreover
correlates positively with the explored node count in every group. The
association of the bound with solver effort is therefore mediated by
the relaxation. Tight capacities give tight root relaxations, and
within the $3{,}600$~s budget the solver rarely moves far from its root
bound. On this benchmark the covering bound annotates relaxation
tightness rather than search difficulty, its direction stable across
groups and its strength moderate within the five-product groups and
weaker pooled. We make no causal claim beyond the controlled
comparison.

\begin{figure}[H]
\centering
\includegraphics[width=0.85\textwidth]{fig_q3_scatter.pdf}
\caption{Terminal gap of \textsc{Mip} under the $3{,}600$~s limit
against the covering bound $k^\ast$, one marker per benchmark
instance, coloured by size group. Within every group larger bounds
accompany smaller gaps, the association Table~\ref{tab:q3groups}
quantifies. The within-group reading is what controls for instance
size.}
\label{fig:q3scatter}
\end{figure}
